\section{Full frames and completion of the coefficient induction}
\label{full:section}

Fix $1\leq t<n<p-1$, and retain the notation
\[
 m=n-t,\qquad G=\GL_m(\Zp),\qquad U_0=P_t\times G.
\]
For compact open subgroups $K\subset P_t$ and $V\subset G$, write
\[
 B_{K,V}=B_{K\times V},\qquad
 B_{\mathrm{conn},V}=\colim_K B_{K,V},\qquad
 B_{\mathrm{full}}=\colim_{K,V}B_{K,V}.
\]
These are algebraic unions.  A cochain with values in one of them uses a
common finite frame level and a common pole bound on its compact source.
The same convention applies when an algebraic union of Frobenius roots is
adjoined.  In particular, none of these unions denotes an analytic
completion.

We complete the induction of Theorem~\ref{thm:coefficient-finiteness}.
Connected-frame cohomology at each fixed matrix level will be finite,
and the perfected full-frame rows will have the norm-character
filtrations needed on subsequent boundary strata.  At the current
pair $(n,t)$ we use compact finiteness and removal of the pole bound,
already supplied by Propositions~\ref{prop:compact-finiteness}
and~\ref{prop:pole-removal}.  No chosen constant classes are needed
for this induction step.

The construction has three stages.  First, a dual coefficient complex
turns the stable transfer image into compact supports on the independent
extension locus.  Scalar symmetry of that support calculation makes
deep translations trivial on cohomology.  Compact-group duality then
gives finiteness at full connected level.  Finally, we follow the actual
root inclusions to determine the norm characters on each perfected
cohomology row.  Section~\ref{full:module-section} will use the same
dual complex to identify the action of the specified constant classes
on the coefficient unit.

All complexes in this section use cohomological grading.  The algebraic
dual $C^*=\Hom_k(C,k)$ has $(C^*)^{-i}=\Hom_k(C^i,k)$ and the usual dual
differential.  For a compact complex we also use its continuous dual
$C^\vee$.  Constant stabilizer cochains act on either dual by the graded
transpose of their cup action.  Equivalently, the dual is a module over
the same cochain algebra in the derived category, with the duality signs.

\subsection{The cutoff after finiteness}

Let $\mathcal E_U(N)$ be the coefficient bundle associated to a finite
translation representation $N$ at a joint frame level $U$ on which its
auxiliary actions are defined.  Fix once and for all a bounded finite
projective resolution for $P_n$ and use it to write
\[
 C_U^\bullet(N)=C^\bullet(P_n,\mathcal E_U(N)).
\]
Its degrees lie in $[0,n^2]$.  The associated-bundle functor is exact on
the coefficient modules under consideration: finite flat descent and
module splittings on the punctured affine base give exactness before
cohomology.  It remains exact with the stipulated continuous parameters.

The regular translation representations will be denoted $N_r$, where
$r$ ranges through the divisible progression of
Lemma~\ref{lem:cartier-transfer}.  Thus $\mathcal E_U(N_r)$ is the
corresponding lift/root algebra.  The transition $N_{r'}\to N_r$ is
normalized Hopf integration.  In stationary coordinates the induced
cochain transition is the transfer operator, denoted $S$ in this
subsection.  The progression is chosen to fix $k$, so this operator is
$k$-linear.

\begin{lemma}[Finite stable images]
\label{full:stable-cutoff}
At a sufficiently deep fixed frame level, let $P=x^{-b}L$ be a
transfer-stable compact lattice as in
Proposition~\ref{prop:transpose-comparison}.  The map
\[
 H^i(P_n,P)\longrightarrow H^i(P_n,B_U)
\]
identifies the stable images of transfer on its source and target.
Consequently the transpose comparison remains valid when the compact
lattice is replaced by the entire bounded-pole coefficient.
The resulting comparison of dual direct-limit complexes preserves
constant stabilizer cochains and finite-level trace maps.
\end{lemma}

\begin{proof}
Put $M=H^i(P_n,P)$ and $X=H^i(P_n,B_U)$.  Both are finite-dimensional:
this is Proposition~\ref{prop:compact-finiteness} for $M$ and
Proposition~\ref{prop:pole-removal} for $X$.  Their comparison kernel is
locally transfer-nilpotent.  Indeed a primitive in the bounded-pole
complex has some finite pole bound, and sufficiently many transfers
carry that primitive into $P$.

Choose cocycle representatives for a basis of $X$.  They have a common
pole bound, which is carried into $P$ by an iterate of transfer.  Hence
$S^aX$ is contained in the image of $M$ for some $a$.  On a finite vector
space, the stable image $X_{\mathrm{inv}}=\bigcap_a S^aX$ is the largest
subspace on which $S$ is invertible; the complementary generalized
zero-eigenspace is nilpotent.  The same holds for $M$.  The two preceding
observations show that $M_{\mathrm{inv}}\to X_{\mathrm{inv}}$ is
surjective and has zero kernel.

The direct limit of the dual transfer system is the dual of the stable
image.  This proves the comparison on cohomology in every degree.  For
the assertion on complexes, use the natural zigzag
\[
 C^\bullet(P_n,B_U)^*
   \longrightarrow C^\bullet(P_n,P)^*
   \longleftarrow C^\bullet(P_n,P)^\vee
\]
and take the direct limits of transpose transfer.  The second arrow is
the inclusion of continuous functionals.  The compact cochain terms and
their closed images have exact continuous duality; algebraic vector
space duality is exact as well.  Since the compact cohomology groups are
finite, the second arrow is a cohomology isomorphism.  The stable-image
argument gives the same assertion for the first arrow after the direct
limit.  Filtered colimits are exact and the complexes are bounded, so
the zigzag consists of quasi-isomorphisms after that limit.

All its arrows are restriction, inclusion of functionals, or transpose
transfer.  Their compatibility with cup products and with finite-frame
trace is the compatibility in
Lemma~\ref{lem:cartier-transfer} and
Proposition~\ref{prop:transpose-comparison}.  No choice of a basis of
$M$ or $X$ is used to define these arrows.
\end{proof}

\subsection{A dual complex with translation action}

Take algebraic duals of the regular-representation complexes and form
the filtered colimit over translation levels and joint frame levels:
\begin{equation}
\label{full:dual-complex}
 \cD=\colim_{U,r} C_U^\bullet(N_r)^*.
\end{equation}
In the translation direction the arrows are dual to Hopf transfer; in
the frame direction they are dual to finite-level trace.  This notation
means a filtered colimit of complexes, and does not require every
displayed arrow to have been written as an inclusion.  Each finite
diagram is taken at one sufficiently deep frame level.  The result is a
bounded complex of rational $T$-representations, with a commuting smooth
$U_0$-action and with the constant stabilizer cochain action.

Here a rational representation of the affine group scheme $T$ is a
comodule for its coordinate Hopf algebra.  Each vector belongs to a
finite-dimensional subcomodule and uses a finite translation level.
This is distinct from an action of the group of geometric points of
$T$.  Write $T_h=[p]^hT$, with $h$ in the same divisible progression as
above.  For $t>1$ its distribution algebra is, after stationary
identification,
\[
 \Lambda_h\simeq k[[D_1,\ldots,D_d]],\qquad d=m(t-1).
\]
For $t=1$, $T_h$ is diagonalizable and its invariants are exact.

\begin{lemma}[Trace descent for the dual complex]
\label{full:trace-descent}
Let $U=K\times V$.  There is a natural equivalence
\begin{equation}
\label{full:dual-descent}
 \RGamma\bigl(U,\RGamma(T_h,\cD)\bigr)
 \simeq C^\bullet(P_n,B_{K,V}^{(h)})^*,
\end{equation}
where $B_{K,V}^{(h)}$ denotes the corresponding lift/root coefficient.
Taking $h=0$ means the trivial translation representation.  The
equivalence is compatible with transpose cup actions.  Under it,
corestriction in $K$ or $V$ is dual to pullback of bounded-pole
coefficients.  Power transport identifies the questions about
$B_{K,V}^{(h)}$ with those about $B_{K,V}$, retaining ordinary
$\Fp$-cochains.
\end{lemma}

The notation $B_{K,V}^{(h)}$ in \eqref{full:dual-descent} is the
$U=K\times V$ instance of the regular-representation coefficient
$\mathcal E_U(N_h)$ described above, at the $h$th lift/root level of
\eqref{rings:root-torsor}.

\begin{proof}
First keep the translation level finite.  If $U'\triangleleft U$ is a
smaller frame level, its coefficient algebra is a finite torsor for
$H=U/U'$.  Over the quotient algebra $A$, faithfully flat
trivialization makes it a regular $A[H]$-module.  Projectivity over
$A[H]$ descends, and $A[H]$ is free as a $k[H]$-module; hence the
additive coefficient module is projective for the finite deck group.
The same argument applies to the descended coefficient bundles.
Coinvariants identify with the lower-level coefficient by trace: on a
trivialized torsor this is the augmentation of the regular module.
The fixed finite projective $P_n$-resolution takes equivariant finite
sums and retracts, so both assertions hold on its cochain terms.

Dual terms are consequently acyclic for deck-group invariants, with
invariants equal to the required dual coefficient at the lower level.
Passing to the smooth union preserves this calculation.  Indeed a
continuous cochain with discrete smooth values has finite image on a
compact source; that image and all its stabilizers are visible at one
finite quotient.  Thus the continuous bar calculation is the filtered
union of these finite calculations.

For translations, test a term of~\eqref{full:dual-complex} against any
finite-dimensional rational representation $W$.  Its equivariant Hom
is the dual of the associated-bundle functor evaluated on $W^*$.
That functor is exact.  Hence the term is injective in the category of
rational $T$-representations.  Here finite tests suffice: in extending
a map from a subcomodule, any new vector lies in a finite-dimensional
subcomodule $W$.  Its intersection with the domain is finite
dimensional, so exactness extends the map across $W$ and then across
their sum.  The maximal-extension argument proves injectivity.  The
restriction to $T_h$ remains injective: induction across the finite
quotient $T/T_h$ is exact, as is seen on its finite locally free Hopf
algebra.

Translation invariants of these terms are therefore computed without
higher termwise cohomology.  They recover the dual coefficient of the
trivial translation representation; $T_h$-invariants recover the
indicated finite lift level.  The bounded complex permits taking this
calculation before the compact frame-group calculation.  This proves
\eqref{full:dual-descent}.

At a finite deck level, dualizing pullback gives trace.  The induced map
on the invariant calculation is precisely corestriction, by the
regular-module norm formula.  This proves the assertion about frame
transitions, including its direction.  The remaining assertions follow
from the equivariant powering identities of
Lemma~\ref{lem:cartier-transfer}.  In particular, powering preserves
constant $\Fp$-cochains even before a scalar extension has been chosen
to make every stationary operator $k$-linear.
\end{proof}

\begin{proposition}[Full-frame support comparison]
\label{prop:full-frame-comparison}
Let $R=(\cO_E/\varpi)^a$, with $\varpi=p^\flat$, and let
$Z=(N_t^m)_{\mathrm{ind}}$ be the independent extension locus of
Theorem~\ref{thm:geometric-quotient}.  Use the geometric coefficient
$\mathscr A=(\cO^{\flat,+}/\varpi)^a$ of
Section~\ref{sec:support}.  There is an equivalence
\begin{equation}
\label{full:eq-three}
 \cD\otimes_k R
 \simeq \RGamma_c(Z;\mathscr A)\otimes\langle\eta\rangle[m+n^2]
\end{equation}
after forgetting translations.  It is $U_0$-equivariant and linear for
the action of ordinary constant stabilizer cochains.  It is also
compatible with all finite-frame trace maps used in
Lemma~\ref{full:trace-descent}.
\end{proposition}

\begin{proof}
Apply Proposition~\ref{prop:transpose-comparison} at a sufficiently
deep fixed frame level.  Lemma~\ref{full:stable-cutoff} replaces its
compact lattice by the bounded-pole coefficient without changing the
dual transfer colimit.  We obtain, naturally in the frame level, the
comparison with the compact support calculation on $[Z/U]$, with the
shift $m+n^2$ and the volume character shown in
\eqref{full:eq-three}.

Now pass up all finite frame covers using transpose trace.  On the
geometric presentation this removes compact invariants.  More
explicitly, finite cover descent computes the supported complexes with
locally constant parameters on the compact deck groups.  A finite
diagram of those parameters is defined at one finite quotient.
Passing to all covers therefore gives the supported complex on $Z$.
This is the same smooth-union descent used in
Lemma~\ref{full:trace-descent}.  The buffered support and coefficient
comparisons of Proposition~\ref{prop:transpose-comparison} are made on
those finite diagrams before taking the union, so their arrows pass to
this limit as well.

Continuous compact-group duality supplies the transpose of the cup
action, the residue pairing supplies the shift by $m$, and the
geometric torsor comparison supplies the structural map to $BP_n$.
These are comparisons of complexes with their actions.  Thus the
result retains the asserted constant-cochain action, rather than only
the dimensions of its cohomology groups.
\end{proof}

\subsection{Finite cohomology and the scalar assertion}

\begin{lemma}[Homotheties and deep translations]
\label{full:scalar-translation}
Let $M$ be a rational $T$-representation with a compatible action of a
scalar $a\in1+p\Zp$, $a\ne1$, conjugating $T$ by multiplication by
$a$ or $a^{-1}$.  If $a$ acts as the identity on $M$, then $T_h$ acts
trivially on $M$ for sufficiently large $h$.  A single $h$ works for a
family of such modules on which this same scalar is the identity.
\end{lemma}

\begin{proof}
Suppose first that $t>1$.  In the distribution algebra $\Lambda$, the
ideal generated by $f\circ[a]-f$ annihilates $M$; replacing $a$ by
$a^{-1}$ does not change the argument.  The formal group identity
\[
 [a-1](z)=F([a](z),[-1](z))
\]
shows that the coordinates of $[a-1]$ belong to this ideal.  This follows
also by setting the two inputs of $F(X,[-1](Y))$ equal: its coordinates
vanish when $X=Y$, and hence lie in the ideal generated by $X_i-Y_i$.
Writing $a-1=p^h u$ with $u\in\Zp^\times$ shows that the ideal contains
the coordinates of $[p]^h$.  The image of the augmentation ideal under
the distribution-algebra map for $T_h\hookrightarrow T$ therefore
annihilates $M$.  This is exactly triviality of the $T_h$-action.

At $t=1$, decompose $M$ into characters of the diagonalizable
translation group.  Conjugation by $a$ carries the character $\chi$ to
$a\chi$.  If $a$ acts identically on $M$, every character occurring in
$M$ satisfies $(a-1)\chi=0$.  It is consequently trivial on $T_h$ for
$h\geq v_p(a-1)$.  Both arguments depend only on $a$, not on a choice
of vectors in $M$.
\end{proof}

We will also use the following compact duality calculation.  The
orientation character of $P_t$ is trivial, since the determinant of its
adjoint action is one.

\begin{lemma}[Corestriction to full connected level]
\label{full:corestriction}
Let $M$ be a finite-dimensional smooth $k$-representation of $P_t$.
For a cofinal sequence of sufficiently small compact open subgroups
$K\subset P_t$, inverse limit with corestriction gives
\[
 \varprojlim_{K,\mathrm{cor}} H^c(K,M)=0\quad(c\ne t^2),
 \qquad
 \varprojlim_{K,\mathrm{cor}} H^{t^2}(K,M)\simeq M.
\]
The last isomorphism uses a fixed orientation and is natural for
commuting endomorphisms of $M$.  More generally, for a bounded smooth
$P_t$-complex $C$ with finite-dimensional cohomology, it gives
\begin{equation}
\label{full:connected-complex-limit}
 \varprojlim_{K,\mathrm{cor}} H^\ell\bigl(\RGamma(K,C)\bigr)
       \simeq H^{\ell-t^2}(C).
\end{equation}
This identification is natural for maps of coefficient complexes,
including degree-shifted maps.  Normal connected levels retain the
commuting actions and the action of $P_t$ by conjugation.
\end{lemma}

\begin{proof}
Shrink $K$ so that it acts trivially on $M$.  Compact group duality
identifies the dual of corestriction in degree $c$ with restriction in
degree $t^2-c$, with the dual coefficient and the orientation character.
The latter character is trivial here.  In a restriction colimit over
open subgroups, positive-degree continuous cohomology with finite
discrete coefficients vanishes: a normalized cocycle and its finite
diagram factor through a finite quotient, and restriction to the
kernel makes its positive-degree class zero.  Degree zero retains the
underlying coefficient.  This proves the assertion by duality.

For the complex assertion, use the bounded spectral sequence
\[
 E_2^{c,j}(K)=H^c(K,H^j(C))
       \Longrightarrow H^{c+j}\bigl(\RGamma(K,C)\bigr).
\]
All its groups at a fixed level are finite over the finite field $k$,
so their inverse systems are Mittag--Leffler.  Taking the inverse limit
is exact on these systems and retains only the column $c=t^2$.
There are finitely many groups and pages in each total degree; no
derived inverse-limit term remains.  This proves
\eqref{full:connected-complex-limit}.  Naturality of the spectral
sequence and of compact duality proves naturality also for a map
$C\to C'[r]$.  Conjugation preserves the chosen orientation because
its adjoint determinant is one, so normal levels preserve the stated
$P_t$-action as well.
\end{proof}

\begin{proposition}[Completion of the joint induction]
\label{full:joint-completion}
For every compact open $V\subset G$, the graded vector space
\[
 H^*(P_n,B_{\mathrm{conn},V})
\]
has finite total dimension.  There is an integer $e$, depending on
$(n,t)$, such that the scalar subgroup $1+p^e\Zp\subset G$ acts
identically on $H^*(P_n,B_{\mathrm{full}})$.  The same scalar subgroup
acts identically after adjoining the algebraic union of Frobenius
roots.  Together with Proposition~\ref{full:perfected-row-characters}
below, these statements complete the full-frame part of
Theorem~\ref{thm:coefficient-finiteness} at the current induction step.
\end{proposition}

\begin{proof}
By Proposition~\ref{prop:compact-support}, every scalar in
$1+p\Zp$ acts identically on the compact support cohomology of $Z$:
it fixes the plane and flag parameters and has trivial reduced
orientation character.  Its character on the volume line is also
trivial.  Proposition~\ref{prop:full-frame-comparison} and faithful
almost scalar extension therefore give the same scalar identity on
$H^*(\cD)$.

Choose such an $a\ne1$.  By
Lemma~\ref{full:scalar-translation}, a single $T_h$ acts trivially on
all these cohomology groups.  Increase $h$ into the prescribed
divisible progression.  The translation hypercohomology spectral
sequence then has rows
\begin{equation}
\label{full:translation-rows}
 H^q(\cD)\otimes_k
 \bigwedge^b(\mathfrak m_{\Lambda_h}/\mathfrak m_{\Lambda_h}^2)^*,
 \qquad 0\leq b\leq d.
\end{equation}
This is the augmentation Koszul calculation for the regular
distribution algebra.  The induced auxiliary action on Ext is the
linear action on its cotangent space, even though an auxiliary
automorphism of the full distribution algebra may be nonlinear.  At
$t=1$ only the column $b=0$ is present, by exactness of invariants.

Apply $V$-cohomology to the rows in
\eqref{full:translation-rows}.  After tensoring with $R$, the flag
resolutions of Proposition~\ref{prop:compact-support} and Shapiro's
lemma reduce their cohomology to that of compact parabolic subgroups
with finite-dimensional coefficient representations.  Those groups
are compact analytic groups without $p$-torsion; their finite
projective resolutions give finite-dimensional cohomology in a
bounded range.  There are only finitely many flag types and finite
exterior factors.  The resulting bounds descend before almost scalar
extension by the bounded almost-length argument of
Corollary~\ref{support:finite-invariants}: every finite-dimensional
subspace of a $k$-cohomology group contributes its dimension to the
length after tensoring with $R$, whereas the finite flag and compact
group resolutions give one upper bound for that length.
Smooth cohomology commutes with this flat tensor, by the fixed finite
projective group resolution.  Thus the $V$-cohomology groups just
obtained are finite-dimensional smooth $P_t$-modules.

Take $K$-cohomology next and then inverse limit with corestrictions in
$K$.  Lemma~\ref{full:corestriction} retains only the top connected
group degree $t^2$, with the underlying finite-dimensional
coefficients.  All cohomological ranges are bounded by the dimensions
of the compact analytic groups, the support range, and $d$.
Mittag--Leffler therefore permits using the inverse limit on these
spectral sequences.  By Lemma~\ref{full:trace-descent}, the result is
the algebraic dual of
\[
 \colim_K H^*(P_n,B_{K,V}^{(h)}).
\]
It is finite-dimensional, so the displayed colimit is
finite-dimensional as well.  Power transport identifies it with the
asserted group for $B_{\mathrm{conn},V}$.  Rectangular subgroups form a
cofinal family, so this proves the finiteness statement at every
compact open $V$.

It remains to check that a scalar subgroup can be chosen uniformly
when both frame levels vary.  The scalar $a$ above is central in $G$,
so it acts trivially by conjugation on every $V$.  On
\eqref{full:translation-rows} it is the identity: it is the identity
on $H^q(\cD)$, and its action on the cotangent factor is given by
reduction of scalar multiplication, hence is also the identity.
It therefore acts identically on the successive spectral-sequence
terms used above.  The filtration lengths are bounded independently
of $K,V$: the dimensions are always $t^2$ and $m^2$, the translation
range is $[0,d]$, and the complex $\cD$ is bounded by $n^2$.
A further fixed $p$-power of $a$ consequently acts identically on the
cohomology of every sufficiently deep rectangular-level calculation:
an operator acting identically on a filtration of length $l$ satisfies
$(a-1)^l=0$, and $a^{p^j}-1=(a-1)^{p^j}$ in characteristic $p$.
These deep rectangular levels are cofinal in the full-frame union.

Dualizing and passing to the algebraic frame union proves scalar
identity on $H^*(P_n,B_{\mathrm{full}})$.  Its action is smooth, so
identity for a topological generator implies identity for the closed
open scalar subgroup it generates.  Finally, every root stage is
identified by powering with the same finite-level calculation, with
the same scalar action.  Cohomology commutes with the algebraic root
union under the common-stage convention, so the same subgroup works
there.  This proves the scalar assertion used on strata of higher
starting height in the subsequent induction steps.
\end{proof}

\subsection{Characters after perfecting the full frame}

The boundary induction needs a norm-character filtration on
each individual matrix-group cohomology row of the perfected
full-frame coefficient. To obtain it, we dualize the actual root
inclusions: their translation corestrictions remove the lower
translation degrees in the root limit, and the remaining top
translation line combines with the invariant volume to give
a norm character.

In this subsection the starting height is denoted by $s$.
Fix an induction pair $(n,s)$, with
\[
 1\leq s<n<p-1,\qquad m=n-s,\qquad
 G=\GL_m(\Zp),\qquad d=m(s-1).
\]
For $s>1$, put
\[
 E_h^b=\bigwedge^b
       (\mathfrak m_{\Lambda_h}/\mathfrak m_{\Lambda_h}^2)^*,
       \qquad 0\leq b\leq d.
\]
For $s=1$, put $E_h^0=k$ and retain only this degree.
We use the compact and bounded-pole assertions already proved at this
pair, the finiteness and scalar assertions of
Proposition~\ref{full:joint-completion}, and the full-frame comparison.
The additional assertion proved here is established at the end of that
same induction step.  In a subsequent boundary step $(n,t)$ it is used
only with $s>t$.

Choose the finite field $k$ to contain $\F_{p^s}$, as in the fixed Honda
conventions.  Let
\[
 P_s^1=1+\mathfrak p_{D_s},\qquad
 \Delta_s=\mu_{p^s-1}\subset P_s,\qquad
 \nu_s:P_s\xrightarrow{\operatorname{Nrd}}\Zp^\times
                         \longrightarrow\Fp^\times.
\]
Reduction identifies $\Delta_s$ with $\F_{p^s}^{\times}$, and
$\nu_s$ restricts to the finite-field norm on $\Delta_s$.
Write $\Delta_s^0=\ker(\nu_s|_{\Delta_s})$.
A $k[\Delta_s]$-module has only norm characters if
$\Delta_s^0$ acts identically on it.  Since
$|\Delta_s|$ is prime to $p$, this condition is preserved by
subobjects, quotients, extensions, and filtered colimits.
It is also detected after the faithful almost scalar extension
$k\to R$ used in the support comparison.

Use the same divisible progression of integers $h$ as in
Lemma~\ref{lem:cartier-transfer}; in particular its powering
isomorphisms fix $k$.  Put
\[
 B_{K,V}^{\mathrm{perf}}=\colim_h B_{K,V}^{(h)},\qquad
 B_{\mathrm{conn},V}^{\mathrm{perf}}
     =\colim_{K,h}B_{K,V}^{(h)},\qquad
 B_{\mathrm{full}}^{\mathrm{perf}}
     =\colim_{K,V,h}B_{K,V}^{(h)} .
\]
Every colimit is algebraic, with the common-stage and common-pole
convention for continuous $P_n$-cochains.  The root transition is the
actual inclusion into the next root ring.  Its direction differs from
the transfer used to construct $\cD$.

The four maps below keep these constructions distinct.  Write
$C_U^{(h)}=C^\bullet(P_n,B_U^{(h)})$ and let $U'\subset U$ be a
finer frame level; $r'>r$ and $h'>h$ denote deeper translation and
root levels.  The dual identifications in the last two rows
are those of Lemma~\ref{full:trace-descent} and
Lemma~\ref{full:root-corestriction} below.
\begin{center}
\renewcommand{\arraystretch}{1.2}
\begin{tabular}{@{}p{.17\textwidth}p{.37\textwidth}p{.37\textwidth}@{}}
\toprule
Map & Coefficient direction & Transpose and use\\
\midrule
Hopf transfer &
$C_U^\bullet(N_{r'})\to C_U^\bullet(N_r)$ &
Direct transition in the translation index defining $\cD$.\\
Frame trace &
$C_{U'}^\bullet(N_r)\to C_U^\bullet(N_r)$ &
Direct transition in the frame index defining $\cD$; geometric pullback.\\
Root inclusion &
$C_U^{(h)}\to C_U^{(h')}$ &
Corestriction from $T_{h'}$ to $T_h$; dualizes perfection.\\
Frame pullback &
$C_U^{(h)}\to C_{U'}^{(h)}$ &
Corestriction from $U'$ to $U$; dualizes the connected-frame union.\\
\bottomrule
\end{tabular}
\end{center}
Stationary powering identifies individual coefficient problems; it is
none of these transition maps.  The next lemma computes the third row,
which is the additional input needed to control characters after
perfection.

\begin{lemma}[Root inclusions and translation corestriction]
\label{full:root-corestriction}
In the equivalence of Lemma~\ref{full:trace-descent}, the transpose of
the inclusion
$B_{K,V}^{(h)}\hookrightarrow B_{K,V}^{(h')}$
for $h'>h$ is finite-Hopf corestriction
\[
 \RGamma(T_{h'},\cD)\longrightarrow\RGamma(T_h,\cD).
\]
This identification commutes with the compact frame-group
corestrictions and with the action of $P_s$ by conjugation on
normal connected levels.

Suppose $s>1$, and suppose $T_h$ acts trivially on every $H^j(\cD)$.
On a cohomology row the displayed transition is zero in
translation degrees $b<d$ and is the oriented identity in degree
$d$, up to a nonzero normalization scalar.
\end{lemma}

\begin{proof}
For $s=1$, the finite diagonalizable translation representations
split into character summands.  The transpose of inclusion is the
projection dual to inclusion of those summands, which is normalized
finite-Hopf corestriction; it preserves the trivial character.
This proves the first assertion in that case.

Now suppose $s>1$.  At a finite translation level use the
invariant trace pairing of
Lemma~\ref{lem:cartier-transfer} to identify the regular
representation with its distribution-algebra model
\[
 N_q=k[D_1,\ldots,D_d]/(D_1^q,\ldots,D_d^q).
\]
The constant function corresponds to the normalized top socle
element.  Inclusion of functions at the next finite quotient is
therefore, in this model,
\[
 N_q\longrightarrow N_{q'},\qquad
 f\longmapsto
       \left(\prod_iD_i^{q'-q}\right)f .
\]
Here compatible invariant pairings give the displayed formula;
for other Frobenius-pairing coordinates it is multiplied by
a unit of the target algebra.
Its image is the submodule annihilated by all $D_i^q$.
The opposite-direction normalized Hopf transfer is the quotient
$N_{q'}\twoheadrightarrow N_q$.
These descriptions follow either from the invariant pairing or
by pairing the two maps with the monomial basis; the normalization
fixes the top functional.  The root-ring presentation in
Proposition~\ref{prop:frame-torsor} identifies the first map with
the actual root inclusion.  Taking its dual in the regular-module
construction of $\cD$ gives corestriction from $T_{h'}$ to $T_h$.
Thus the transition here is distinct from the transpose of Hopf
transfer used to define $\cD$.

Here is the cochain calculation, with directions explicit.
Write
\[
 A=\Lambda_h=k[[x_1,\ldots,x_d]],\qquad
 A'=\Lambda_{h'}=k[[y_1,\ldots,y_d]],\qquad
 y_i\longmapsto x_i^q ,
\]
where $q>1$ is the relative exponent in the divisible progression.
The Koszul complex $K_A(x_1,\ldots,x_d)$ resolves $k$, and
\[
 A\otimes_{A'}K_{A'}(y_1,\ldots,y_d)
       =K_A(x_1^q,\ldots,x_d^q)
\]
resolves $A/(x_1^q,\ldots,x_d^q)$.
The socle inclusion $k\to A/(x_i^q)$ is lifted by the chain map
\begin{equation}
\label{full:root-koszul-map}
 e_I\longmapsto
       \left(\prod_{j\notin I}x_j^{q-1}\right)f_I .
\end{equation}
Indeed both sides of its differential identity have, for an
index $i\in I$, the factor
$x_i^q\prod_{j\notin I}x_j^{q-1}$.
Applying $\Hom_A(-,\cD)$ gives
\[
 \RHom_{A'}(k,\cD)\longrightarrow\RHom_A(k,\cD),
\]
the asserted corestriction.  On a row $M=H^j(\cD)$ annihilated
by $(x_1,\ldots,x_d)$, its degree-$b$ map is multiplication by
the complementary product in \eqref{full:root-koszul-map}.
It is zero if $b<d$ and the identity if $b=d$.

Multiplying the socle inclusion by a unit multiplies these
components by that unit; on $M$ its effect is its nonzero
augmentation.  Thus the vanishing and top isomorphism do
not require a choice of compatible unit normalizations.
The maps themselves are the finite Hopf inclusion and its
transpose, so they are natural under all frame actions,
including nonlinear augmented automorphisms of the distribution
algebra.  The cohomology-row calculation is therefore
equivariant.  We do not need to identify the full complexes
at a fixed root level with their top rows.
All finite coefficient diagrams commute with frame changes and
their traces, giving the remaining compatibilities.
\end{proof}

\begin{lemma}[The perfected connected-frame calculation]
\label{full:perfected-connected-calculation}
Choose $h_0$ so that $T_{h_0}$ acts trivially on $H^*(\cD)$,
and set $A_h=\RGamma(V,\RGamma(T_h,\cD))$.
For every compact open $V\subset G$, there are natural
$P_s$-equivariant identifications
\begin{equation}
\label{full:perfected-connected-dual}
 \Hom_k\!\left(
     H^i(P_n,B_{\mathrm{conn},V}^{\mathrm{perf}}),k\right)
 \ \simeq\
 \varprojlim_{h,\operatorname{cor}} H^{-i-s^2}(A_h).
\end{equation}
All vector spaces in this formula are finite dimensional.
Let $\chi$ be a character of $\Delta_s$.
If
\begin{equation}
\label{full:top-row-character-test}
 H^a(V,H^q(\cD)\otimes E_h^d)_\chi=0
       \quad\text{for every }a,q\text{ and }h\geq h_0,
\end{equation}
then the $\chi$-isotypic part of the dual in
\eqref{full:perfected-connected-dual} is zero.
For $s=1$, put $d=0$ and $E_h^0=k$.
\end{lemma}

\begin{proof}
The finite flag calculation and the translation-row spectral
sequence in the proof of
Proposition~\ref{full:joint-completion} show that $H^*(A_h)$
has finite total dimension.
This uses only $V$-cohomology of the support flag representations;
no cohomology of a coefficient twisted by an arbitrary
$P_s$-representation is used.
Finiteness is obtained before almost scalar extension by
Corollary~\ref{support:finite-invariants}.
The dimensions of the cohomological ranges are bounded by
$n^2$, $m^2$, and $d$, independently of $h$.

First take the connected-frame limit at a fixed $h$.
Choose a cofinal sequence of normal principal open subgroups
$K\triangleleft P_s$.
Lemma~\ref{full:trace-descent} identifies
$\RGamma(K,A_h)$ with the finite-frame dual.
Compact duality and inverse limit with corestriction therefore give
\begin{equation}
\label{full:connected-limit-before-roots}
 \Hom_k\!\left(
      H^i(P_n,B_{\mathrm{conn},V}^{(h)}),k\right)
       \simeq H^{-i-s^2}(A_h).
\end{equation}
This is the calculation of Lemma~\ref{full:corestriction},
applied to the bounded finite cohomology rows of $A_h$.
It retains the full $P_s$-action: an element of $P_s$ acts
on $K$ by conjugation as well as on its coefficient.
In compact duality its additional character is the determinant
of the adjoint action on $\Lie(P_s)$, which is one.
For example, after a splitting field extension the Lie algebra is
$M_s$ and conjugation has determinant one.
Thus the top-degree survivor in
\eqref{full:connected-limit-before-roots} carries exactly the
original $P_s$ coefficient action.
Normality of $K$ makes these actions compatible throughout the
inverse system.
The finite cohomology groups are Mittag--Leffler, so no derived
inverse-limit term occurs.

All these operations precede scalar extension.
The actual power-transport isomorphism at a fixed $V$ identifies
the dimensions of
$H^*(P_n,B_{\mathrm{conn},V}^{(h)})$
with the same fixed finite dimensions for every $h$ in the
progression.  These isomorphisms give a uniform bound on dimensions;
they do not identify the root inclusions with isomorphisms.
Consequently the root direct limit is finite dimensional, and
duality identifies it with the inverse limit of the finite
spaces in \eqref{full:connected-limit-before-roots}.
Every cohomological inverse system here is Mittag--Leffler.

This proves \eqref{full:perfected-connected-dual}.
It remains to prove the character test.
For $s>1$, fix a row $M=H^q(\cD)$ and consider the
bounded spectral sequence
\[
 H^a(V,H^b(T_h,M))
       \ \Longrightarrow\
 H^{a+b}\bigl(V,\RGamma(T_h,M)\bigr).
\]
Here $H^b(T_h,M)=M\otimes E_h^b$.
By Lemma~\ref{full:root-corestriction}, root corestriction
acts as zero on every row with $b<d$.
On the $\chi$-part the remaining $b=d$ row is zero by
\eqref{full:top-row-character-test}.
Thus every transition is zero on this row page after
projecting to $\chi$.
The ranges $0\leq a\leq m^2$ and $0\leq b\leq d$ give a
uniform finite filtration on its abutment.
A transition zero on the associated graded lowers that
filtration, so a bounded number of successive transitions
is zero on the $\chi$-part of the abutment.

Next use the bounded cohomology filtration of $\cD$ to compute
$A_h$ from these row calculations.
Its number of nonzero cohomology rows is at most $n^2+1$.
Repeating the same filtration argument therefore gives one
integer $N$, independent of $h$, such that the composite of
$N$ successive root transitions is zero on
$H^j(A_h)_\chi$ in every degree.
For example, one may enlarge the bound to
$(n^2+1)(m^2+d+1)$.
This two-stage argument retains all derived translation and
$V$-extensions; it assumes no equivariant formality.
The $\chi$-part of the inverse system is uniformly pro-zero.
Its inverse limit is consequently zero, as is its derived
inverse-limit contribution.

For $s=1$, translation invariants are exact character projection.
Once $T_{h_0}$ acts trivially on $H^*(\cD)$, only $b=0=d$
occurs.  Condition \eqref{full:top-row-character-test} kills
the corresponding $\chi$-part directly by the bounded
cohomology filtration of $\cD$.
Throughout the proof, inverse limits were evaluated on the
finite $k$-cohomology systems before tensoring with $R$;
no interchange of an infinite inverse limit with almost
scalar extension is required.
\end{proof}

\begin{lemma}[The remaining tame characters]
\label{full:perfected-support-characters}
The $P_s$-action on
$H^i(P_n,B_{\mathrm{conn},V}^{\mathrm{perf}})$
has only norm characters on $\Delta_s$, for every $i$ and every
compact open $V\subset G$.
\end{lemma}

\begin{proof}
We record the connected-frame characters in the support calculation.
The orientation line for a rational plane of dimension $r$ comes
by scalar extension from a one-dimensional $\Fp$-space, naturally
for automorphisms of $N_s$.
Indeed the natural compact primitive comparison
\[
 \RGamma_c(N_s^r,\Fp)\otimes_{\Fp}R
       \longrightarrow \RGamma_c(N_s^r,\mathscr A)
\]
is computed by the same buffered affine exhaustion and lattice
corestriction calculation as
Lemma~\ref{support:translation-quotient}.
The affine finite-coefficient complex is $\Fp(-r)[-2r]$;
the lattice corestriction limit retains its top
$sr$-dimensional orientation.
Thus the source has one $\Fp$-cohomology line in degree
$(s+2)r$, and the comparison is an equivalence.
The calculation is natural, so it applies also to $P_s$
automorphisms which do not preserve a chosen affine/lattice
presentation.
The compact-parameter construction gives a continuous action
on that finite line.  Its $P_s$ character consequently takes
values in $\Fp^\times$.

Every continuous character $P_s\to\Fp^\times$ is a power of
$\nu_s$.  Its restriction to the pro-$p$ group $P_s^1$ is
trivial; the quotient is the cyclic group
$\F_{p^s}^\times$, and every homomorphism from that group to
$\Fp^\times$ is a power of the finite-field norm.
Denote the resulting orientation character for an $r$-plane
by $\nu_s^{c_r}$; its exact exponent is not needed.
The rational plane and flag parameters are fixed by $P_s$,
since $P_s$ commutes with the action on the $m$ row indices.
Hence each finite flag term for the support row is a
representation with this scalar $P_s$ character.

There are also the volume and top translation Ext lines.
Let $a\in\F_{p^s}^\times$ be a Teichm\"uller element, acting on
the tangent line of the fixed Honda group by $a$.
With the inverse-substitution convention of
\eqref{rings:auxiliary-action}, the characters are
\[
 \langle\eta\rangle(a)=a^{-m},\qquad
 E_{h_0}^d(a)=a^{-m(p+p^2+\cdots+p^{s-1})}.
\]
For the second identity, let $h=[a]$ act on the Honda
group.  Its action on functions is
$h\cdot f=(h^{-1})^*f$, so the dual distribution action is
$\lambda\mapsto(f\mapsto\lambda(h^*f))$.
Under Cartier duality this is pullback by
$h^D(\rho)=\rho\circ h$.
To compute these characters, let
$O=W(\F_{p^s})\subset\operatorname{End}(\Gamma_s)$ act on the
covariant Dieudonne module $M$ over $W(k)$.  The Honda endomorphism
calculation identifies a Teichmuller element $a$ with the
endomorphism $aX$
\cite[Lemma~A2.2.16 and Theorem~A2.2.18]{RavenelGreenBook}.
The $s$ idempotents of $W(k)\otimes_{\Zp}O$ split $M$ into its
eigensummands.  Its semilinear Frobenius commutes with $O$ and is
invertible after inverting $p$, so it permutes the rationalized
eigensummands cyclically and forces their ranks to be equal
\cite[Example~14 and Proposition~15]{ZinkDisplays}.
The covariant Hodge sequence for $M/pM$ has total rank $s$, the sum of the
dual dimension $s-1$ and the dimension one of $\Gamma_s$
\cite[Section~5]{Lau2010Tate}.  Each eigensummand therefore has
rank one, and the residue eigencharacters are
$a,a^p,\ldots,a^{p^{s-1}}$.
In this functorial Hodge sequence the quotient
$\Lie(\Gamma_s)$ has character $a$, and the subspace
$\omega_{\Gamma_s^D}$ is the cotangent space of the
distribution formal group, with characters
$a^p,\ldots,a^{p^{s-1}}$.
Degree-one translation Ext is the linear dual of this cotangent
space; top Ext is its top exterior power.  This gives the stated
negative exponents.
In particular the action used here is cotangent
pullback by $h^D$, not the tangent action of
$\rho\mapsto\rho\circ h^{-1}$ on dual points.
The divisible progression fixes these
characters when transported to $T_{h_0}$.
Therefore
\begin{equation}
\label{full:connected-top-character}
 (\langle\eta\rangle\otimes E_{h_0}^d)(a)
   =a^{-m(1+p+\cdots+p^{s-1})}
   =\nu_s(a)^{-m}.
\end{equation}
For $s=1$ the second line is trivial and the same formula holds.
This is a calculation on the fixed Honda group and its actual
translation representations.  It does not identify a
parameter-dependent universal Lie line with the constant
extension of its special-fibre character.

To apply \eqref{full:top-row-character-test}, use
Proposition~\ref{prop:full-frame-comparison} on each
$H^q(\cD)\otimes E_h^d$ after tensoring with $R$.
By the finite flag resolution, the resulting support rows
have only norm characters of $\Delta_s$, namely the products
of $\nu_s^{c_r}$ with \eqref{full:connected-top-character}.
$V$ commutes with $\Delta_s$.
Its finite projective resolution therefore preserves this
condition, and commutes with the flat scalar extension.
Faithfulness detects
\eqref{full:top-row-character-test} before scalar extension
for every nonnorm $\chi$.
Lemma~\ref{full:perfected-connected-calculation} now kills
all these isotypic pieces of the perfected dual.
The set of norm characters is closed under inversion, so
dualizing proves the assertion on the primal coefficient.
\end{proof}

\begin{proposition}[Norm characters on the perfected individual rows]
\label{full:perfected-row-characters}
At the current induction pair $(n,s)$, for every compact open
$V\subset\GL_{n-s}(\Zp)$ and all $a,i$, the vector space
\[
 Q=H^a\!\left(V,
           H^i(P_n,B_{\mathrm{full}}^{\mathrm{perf}})\right)
\]
is finite dimensional and has a finite $P_s$-stable filtration
whose successive quotients are one-dimensional characters
$\nu_s^j$.
Its dual has a filtration by the inverse characters.
This assertion concerns the representations arising from
these perfected full-frame rows.
\end{proposition}

\begin{proof}
First take the algebraic union over the matrix levels in
Lemma~\ref{full:perfected-support-characters}.
The fixed finite $P_n$-resolution and the common-stage
cochain convention give
\[
 H^i(P_n,B_{\mathrm{full}}^{\mathrm{perf}})
    =\colim_{V'} H^i(P_n,B_{\mathrm{conn},V'}^{\mathrm{perf}}).
\]
The maps here are actual coefficient pullbacks and are
$P_s$-equivariant.
Since $\Delta_s^0$ acts identically on every term, it acts
identically on this individual $P_n$-cohomology group.
This step uses an algebraic coefficient colimit; it does
not deduce an individual row from a total-cohomology
spectral-sequence abutment.

The group $V$ commutes with $P_s$, in particular with
$\Delta_s^0$.  A bounded finite projective resolution for
$V$ computes its cohomology by retracts of finite sums of
the displayed coefficient module.  Thus $\Delta_s^0$
also acts identically on $Q$.
Its finite-dimensionality is the individual-row finiteness
assertion of Lemma~\ref{boundary:row-finiteness}.
The hypotheses of that lemma hold for the perfected
coefficient as well: finite-frame descent identifies its
total $V'$-cohomology with the connected-frame calculation,
and the uniform finite root-stage dimensions used in
Lemma~\ref{full:perfected-connected-calculation} give finite
dimensions after the root union, for every open $V'$.
The cohomological ranges stay bounded by the fixed group
dimensions.

Finally the image $\Pi$ of $P_s^1$ in $\GL_k(Q)$ is a finite
normal $p$-group.  Its augmentation ideal $J\subset k[\Pi]$
is nilpotent.  The finite filtration $J^bQ$ is $P_s$-stable,
and its quotients have trivial $P_s^1$ action.
They therefore split over $k$ into characters of
$P_s/P_s^1=\F_{p^s}^\times$.
Exactness of restriction to $\Delta_s$ and of its character
projections shows that each such character is trivial on
$\Delta_s^0$.  Each is consequently a power of $\nu_s$.
Refine the filtration by these one-dimensional summands.
Duality reverses this finite filtration and inverts its
characters, proving the final assertion.
\end{proof}
